\section{Appendix}

\subsection{Computation of Order Statistics}
\label{app:order-stats}
Here we provide details on how the order statistics in Figure~\ref{fig:wasted} (top row) are calculated. 

\paragraph{Acceptance Event}
To track the ordering for the Unified EB sampler, we look at the first
acceptance per position, as the standard entropy bound of $0.1$ results in
positions rarely changing after the first acceptance. Concretely, for a
canvas of $b$ positions $X_1,\dots,X_b$ denoised over $T$ steps, we record
for each position the step at which it is first accepted,
\begin{equation}
  c_i \;=\; \min\bigl\{\, t \in \{1,\dots,T\} \;:\; X_i \text{ is accepted at step } t \,\bigr\},
\end{equation}
setting $c_i = T$ for the rare positions never accepted. Each generated
canvas therefore yields a acceptance-step vector $c \in \{1,\dots,T\}^{b}$. A
purely autoregressive sampler would produce $c_i = i$, so we measure how
closely $c$ tracks the left-to-right index.

\paragraph{Resolving Ties}
Since our goal is to compare the ranking of positions based on acceptance
step, there must be a way to resolve the ties between tokens that are
accepted on the same step. To resolve this, we take the approach of
averaging the order ranking across slots that are tied, assigning each
position the fractional midrank
\begin{equation}
  r_i \;=\; \bigl|\{\,k : c_k < c_i\,\}\bigr|
        \;+\; \tfrac{1}{2}\bigl(1 + \bigl|\{\,k : c_k = c_i\,\}\bigr|\bigr).
\end{equation}
Given the list of acceptance steps $c = [1, 3, 1, 2, 3, 3]$, we first sort them
in ascending order, $[1, 1, 2, 3, 3, 3]$. The positions accepted on the
first step occupy the first and second ranking slots, so both are assigned
the average rank $1.5$; likewise the three positions accepted on the third
step occupy slots four through six and are each assigned $5$. Returning
these ranks to their original positions yields
$r = [1.5,\, 5,\, 1.5,\, 3,\, 5,\, 5]$. Averaging within tie groups keeps
tied positions exchangeable, so the statistic does not depend on an
arbitrary tie-breaking order, and preserves the rank sum
$\sum_i r_i = b(b+1)/2$.

\paragraph{Spearman $\rho$}
Given the vector of rankings, we compare it against the standard
left-to-right rankings $[1, 2, \dots, b-1, b]$ and take the Pearson
correlation coefficient to be the Spearman rank correlation $\rho$. Writing
$\pi_i = i$ for the reference ranking,
\begin{equation}
  \rho \;=\;
  \frac{\sum_{i=1}^{b} (\pi_i - \bar{\pi})(r_i - \bar{r})}
       {\sqrt{\sum_{i=1}^{b} (\pi_i - \bar{\pi})^{2}}\;
        \sqrt{\sum_{i=1}^{b} (r_i - \bar{r})^{2}}},
  \qquad \bar{\pi} = \bar{r} = \frac{b+1}{2}.
\end{equation}
We use this general form rather than the familiar
$1 - 6\sum_i d_i^{2}/\bigl(b(b^{2}-1)\bigr)$, which is valid only when the
ranks are untied.

\paragraph{Kendall $\tau$}
Spearman rank correlation may not reflect the true behavior, as it is
quadratic in terms of the differences per position. Thus we also calculate
the Kendall $\tau$, as it counts inversions and is thus less sensitive to
large differences. More specifically, for every position pair $(X_i, X_j)$
with $i < j$, the pair is \emph{concordant} if $X_i$ was accepted earlier
than $X_j$, \emph{discordant} if $X_j$ was accepted earlier, and
\emph{tied} if both were accepted on the same step. Over all
$n_0 = \binom{b}{2}$ pairs we count
\begin{equation}
  n_c = \bigl|\{(i,j) : i<j,\; c_i < c_j\}\bigr|, \quad
  n_d = \bigl|\{(i,j) : i<j,\; c_i > c_j\}\bigr|, \quad
  n_t = n_0 - n_c - n_d .
\end{equation}
Because the reference ranking is strictly increasing, it contributes no
ties of its own, and every tied pair is tied in acceptance step. We report the
tie-corrected $\tau_b$, which normalises by the pairs whose order is
resolvable at all,
\begin{equation}
  \tau_b \;=\; \frac{n_c - n_d}{\sqrt{n_0\,(n_0 - n_2)}},
  \qquad
  n_2 \;=\; \sum_{s=1}^{T} \binom{m_s}{2},
\end{equation}
where $m_s$ is the number of positions first accepted at step $s$. We
additionally report the pairwise left-to-right agreement rate, which counts
ties as half-agreements,
\begin{equation}
  P_{\mathrm{l2r}} \;=\; \frac{n_c + \tfrac{1}{2} n_t}{n_0},
\end{equation}
equal to $50\%$ under a random acceptance order and $100\%$ under a purely
left-to-right one.

\paragraph{KDE Estimate}
To get the shape of the distribution, we use the default KDE estimator in matplotlib. 

% \begin{figure}
    \centering
    \includegraphics[width=\linewidth]{images/fig6_first_acceptance_position_scatter.pdf}
    \caption{Scatter plot showing the first accept in the trajectory (y axis) against the position within the canvas (x axis) for both tokens that realize their final identity immediately (top row) v.s tokens that fluctuate between the first acceptance and the final canvas (bottom row).}
    \label{fig:acc_vs_pos}
\end{figure}
% \begin{figure}
    \centering
    \includegraphics[width=\linewidth]{images/fig4_first_acceptance_realization.pdf}
    \caption{Histogram of behavior of positions after first acceptance in EB rule. While the majority of the positions stay the same after their first acceptance, a non trivial amount do change from the initially accepted token. This introduces the potential for performance deterioration.}
    \label{fig:placeholder}
\end{figure}
% \begin{figure}
    \centering
    \includegraphics[width=\linewidth]{images/fig5_first_acceptance_timing_histogram.pdf}
    \caption{Histogram of when the first acceptance for tokens tends to occur in the denoising trajectory, organized by tokens that realize their final token identity by the first acceptance (top row) and tokens that change (bottom row). As visible, both tend to overlap very heavily, meaning that there is no clear heuristic to distinguish between whether a position has changed to its final identity or a transient state.}
    \label{fig:acc_vs_timing}
\end{figure}